\section{Criba lineal y menor factor primo}
\whenuse{Muchas consultas de primalidad/factorizaci\'on para valores acotados por
$N$; precomputar una vez y factorizar luego por el menor primo.}

\begin{lstlisting}[style=cpp]
struct LinearSieve {
    vector<int> primes, smallest_prime;

    explicit LinearSieve(int n) : smallest_prime(max(0, n) + 1, 0) {
        for (int value = 2; value <= n; ++value) {
            if (smallest_prime[value] == 0) {
                smallest_prime[value] = value;
                primes.push_back(value);
            }
            for (int prime : primes) {
                if (prime > smallest_prime[value] || 1LL * value * prime > n)
                    break;
                smallest_prime[value * prime] = prime;
            }
        }
    }

    vector<int> factorize(int value) const {
        assert(1 <= value && value < (int)smallest_prime.size());
        vector<int> factors;
        while (value > 1) {
            factors.push_back(smallest_prime[value]);
            value /= smallest_prime[value];
        }
        return factors;
    }
};
\end{lstlisting}

\complexity{preproceso $O(N)$; factorizaci\'on $O(\log n)$; memoria $O(N)$.}

\section{Criba segmentada: primos en un intervalo}
\whenuse{Listar los primos de $[L,R]$ cuando el intervalo es corto y $R$ es
grande. Debe ser viable cribar hasta $\lfloor\sqrt R\rfloor$ y almacenar
$R-L+1$ marcas; por ejemplo, $R\le10^{12}$ y ancho hasta $10^6$.}
\complexity{Con $B=\lfloor\sqrt R\rfloor$ y $W=R-L+1$, tiempo
$O(B\log\log B+\pi(B)+W\log\log R)$ y memoria $O(B+W)$ bits para las marcas,
m\'as $O(k)$ palabras para devolver los $k$ primos. Las cotas se entienden
asint\'oticamente; intervalos con $R<2$ se resuelven directamente.}

\noindent\begin{minipage}{\linewidth}
\begin{lstlisting}[style=cpp]
// Requiere 0 <= L <= R y memoria/tiempo suficientes para sqrt(R) y el ancho.
vector<ll> segmented_primes(ll L, ll R) {
    assert(0 <= L && L <= R);
    if (R < 2) return {};

    ll root = (ll)sqrtl((long double)R);
    while (root > R / root) --root;
    while (root + 1 <= R / (root + 1)) ++root;

    vector<bool> base((size_t)root + 1, true);
    base[0] = base[1] = false;
    vector<bool> segment((size_t)(R - L) + 1, true);

    for (ll p = 2; p <= root; ++p) {
        if (!base[p]) continue;
        if (p <= root / p)
            for (ll x = p * p; x <= root; x += p)
                base[x] = false;

        // Primer multiplo >= L, sin usar (L+p-1) ni una suma fuera de R.
        ll delta = (p - L % p) % p;
        if (delta > R - L) continue;
        ll start = max(p * p, L + delta);
        if (start > R) continue;
        for (ll x = start; ; ) {
            segment[x - L] = false;
            if (R - x < p) break;
            x += p;
        }
    }

    if (L == 0) segment[0] = false;
    if (L <= 1) segment[1 - L] = false;
    vector<ll> answer;
    for (size_t i = 0; i < segment.size(); ++i)
        if (segment[i]) answer.push_back(L + (ll)i);
    return answer;
}
\end{lstlisting}
\end{minipage}

Ejemplo: \texttt{segmented\_primes(100,120)} devuelve
$101,103,107,109,113$. Se marca desde
$\max(p^2,\lceil L/p\rceil p)$ para no borrar al propio primo $p$.
No construye la criba completa hasta $R$. Tampoco evita el costo de cribar
hasta $\sqrt R$: para un solo entero grande, usar Miller--Rabin.
Si hay muchas consultas con el mismo l\'imite, conviene reutilizar los primos
base; esta versi\'on los recalcula para mantener una interfaz autocontenida.

\section{Factorizaci\'on por divisi\'on por tentativa}
\whenuse{Una o pocas consultas con $n$ moderado, por ejemplo hasta $10^9$.
No necesita criba ni generador aleatorio. Para muchas consultas, usar SPF;
para factores grandes de 64 bits, Miller--Rabin y Brent.}
\complexity{$O(\sqrt n)$ tiempo en el peor caso; $O(\log n)$ memoria de salida.}

\begin{lstlisting}[style=cpp]
vector<ll> trial_factorize(ll n) {
    assert(n >= 1);
    vector<ll> out;
    while (n % 2 == 0) { out.push_back(2); n /= 2; }
    for (ll d = 3; d <= n / d; d += 2)
        while (n % d == 0) { out.push_back(d); n /= d; }
    if (n > 1) out.push_back(n);
    return out;
}
\end{lstlisting}

Devuelve factores ordenados con repeticiones; para $1$, un vector vac\'io.
La condici\'on \texttt{d <= n/d} evita overflow de \texttt{d*d}.

\section{Listar divisores}
\whenuse{Listar todos los divisores de un n\'umero aislado moderado sin necesitar
preprocesamiento.}

\begin{lstlisting}[style=cpp]
vector<ll> divisors(ll n) {
    assert(n > 0);
    vector<ll> answer;
    for (ll d = 1; d <= n / d; ++d) if (n % d == 0) {
        answer.push_back(d);
        if (d != n / d) answer.push_back(n / d);
    }
    sort(answer.begin(), answer.end());
    return answer;
}
\end{lstlisting}

\complexity{$O(\sqrt n)$ tiempo y $O(\tau(n))$ memoria de salida.}
